\chapter{REVISÃO BIBLIOGRÁFICA}
\label{cap1Revisao}
Este capítulo apresenta os principais fundamentos teóricos que embasam o desenvolvimento do ERP proposto para otimização dos processos em marmorarias. São discutidos estudos e conceitos relacionados a ERP, transformação digital em pequenas empresas, automação de processos produtivos, práticas modernas de desenvolvimento de software, computação em nuvem e segurança de aplicações. Ao contextualizar essas áreas, busca-se demonstrar como a literatura sustenta a relevância e a necessidade da solução desenvolvida ao longo deste trabalho.

\section{Sistemas Integrados de Gestão (ERP) e sua Evolução}
Os Sistemas Integrados de Gestão Empresarial (ERP) surgiram como resposta à crescente necessidade de unificação dos processos organizacionais em plataformas únicas e centralizadas. As raízes desses sistemas remontam a 1913, quando o engenheiro Ford Whitman Harris desenvolveu o modelo EOQ (Economic Order Quantity), precursor matemático do controle de estoques. Entretanto, o salto para a automatização ocorreu em 1964, quando a Black e Decker, em parceria com a IBM, implementou o primeiro sistema MRP (Material Requirements Planning) para prever necessidades de materiais baseadas na demanda.
Na década de 1980, o modelo evoluiu para o MRP II (Manufacturing Resource Planning), introduzido por Oliver Wight, que passou a integrar não apenas materiais, mas também recursos humanos, financeiros e capacidade de produção. O termo ERP foi finalmente cunhado pelo Gartner Group no início da década de 1990, marcando a expansão definitiva desses sistemas para todas as áreas da empresa de forma integrada. Atualmente, o mercado global de ERP é um pilar estratégico da economia, avaliado em 65,25 bilhões em 2024 com projeção de atingir 103,95 bilhões até 2029, refletindo uma busca contínua por eficiência operacional e inteligência de dados. A literatura enfatiza que a adoção de ERPs traz melhorias diretas na produtividade e na consistência das informações. Contudo, muitos autores destacam que sistemas genéricos frequentemente não atendem às necessidades de nichos industriais específicos, como o setor de marmorarias, que demanda cálculos complexos de corte e gestão de chapas, tornando os chamados ERPs Verticais (especializados por setor) a solução mais adequada.

\section{Transformação Digital e Automação em Pequenas Empresas}
A transformação digital tem remodelado a forma como empresas de todos os portes estruturam seus processos internos. De acordo com a Gartner, a migração para a computação em nuvem consolidou-se como uma das principais tendências do mercado, sendo que, até 2022, cerca de 30 por cento dos novos sistemas financeiros destinados a empresas de médio porte já utilizavam soluções baseadas em nuvem. Essa democratização tecnológica, impulsionada pelo modelo SaaS, permite que pequenas empresas tenham acesso a ferramentas robustas sem a necessidade de investir em infraestruturas locais caras e complexas.

Entretanto, esse avanço também traz novos desafios. Segundo o relatório Data Breach Investigations Report (DBIR) da Verizon (2020), aproximadamente 43 por cento dos ataques cibernéticos têm como alvo pequenas empresas, evidenciando que a segurança da informação deve ser considerada um requisito essencial em qualquer processo de digitalização.

No segmento das marmorarias, a dependência de planilhas eletrônicas, anotações manuais e processos pouco padronizados ainda é uma realidade. Esse cenário favorece a ocorrência de erros operacionais, retrabalho e dificuldades no controle financeiro e produtivo, tornando a automação um importante diferencial competitivo para empresas do setor.

Estudos demonstram que a digitalização de processos em setores tradicionais pode reduzir em até 40 por cento o retrabalho operacional, além de aumentar significativamente a precisão das medições, dos orçamentos e do controle dos custos de produção. Nesse contexto, o desenvolvimento de um ERP especializado para marmorarias representa não apenas um avanço tecnológico, mas uma necessidade estratégica, alinhada à crescente transformação digital observada em segmentos correlatos, como a marcenaria e a construção civil. A adoção de uma solução específica para esse mercado possibilita maior eficiência operacional, melhor gestão dos recursos e maior competitividade frente às exigências do cenário atual.

\subsection{Da Gestão em Papel à Automatização: Um Exame dos Resultados Práticos}

A transição de um modelo de gestão baseado em papel para um ERP especializado representa um redesenho operacional com impactos quantificáveis na eficiência. No setor de rochas ornamentais e marmorarias, o controle rigoroso do estoque de chapas e o aproveitamento de retalhos são gargalos históricos. De acordo com a literatura de engenharia de produção voltada ao mapeamento de fluxo de valor (MFV), a falta de rastreabilidade digital de sobras de produção é um dos principais fatores geradores de desperdício de matéria-prima, uma vez que retalhos utilizáveis acabam sendo descartados por falta de catalogação sistemática.

Além do desperdício de material, a gestão analógica gera um custo oculto de produtividade. Um estudo global da consultoria McKinsey e Company aponta que a automação e a digitalização de processos em setores manufatureiros tradicionais podem liberar até 20 por cento da capacidade produtiva das equipes, reduzindo drasticamente o tempo gasto em atividades puramente burocráticas - como a digitação manual de orçamentos e a conferência física de estoques. Com a implementação de um ERP vertical, esse tempo é revertido para o planejamento estratégico e a otimização do plano de corte. A centralização dos dados elimina os erros de comunicação entre a medição na obra, o setor de vendas e a linha de produção, garantindo que o produto final corresponda exatamente ao projetado e mitigando os custos de retrabalho.

\section{Arquitetura de Software Moderna e Desenvolvimento de Sistemas Web}

O desenvolvimento de sistemas modernos tem sido fortemente influenciado por abordagens arquiteturais flexíveis, modulares e desacopladas. O avanço das demandas de negócios e a necessidade de alta disponibilidade transformaram a engenharia de software, impulsionando a transição dos sistemas monolíticos tradicionais para ecossistemas baseados em microsserviços, APIs REST e conteinerização. Conforme apontam \citeonline{fowler2014}, a arquitetura de microsserviços decompõe uma aplicação em um conjunto de serviços autônomos, organizados em torno de capacidades de negócio específicas e que se comunicam por meio de protocolos leves. Essa abordagem mitiga o risco de acoplamento excessivo, permitindo escalabilidade granular, ciclo de vida independente para cada módulo e atualizações contínuas sem a interrupção do sistema como um todo, como bem descreve \citeonline{newman2015}.

Nesse cenário de descentralização, a comunicação eficiente entre os componentes torna-se um pilar crítico. O padrão arquitetural REST, introduzido historicamente por \citeonline{fielding2000}, consolidou-se como o modelo predominante para a construção de APIs na Web devido à sua simplicidade, escalabilidade e respeito às restrições do protocolo HTTP, como o estado \textit{stateless} (sem estado).

No contexto do ERP desenvolvido neste trabalho, a arquitetura adotada é a de um monólito modular organizado em camadas (domínio, casos de uso, entrega HTTP e infraestrutura), inspirada na \textit{Clean Architecture}, e não uma decomposição em microsserviços independentes. Ainda assim, princípios como APIs REST \textit{stateless} e autenticação baseada em JWT foram aplicados. A adoção do JWT segue as melhores práticas de segurança para arquiteturas modernas, uma vez que elimina a dependência de sessões mantidas no servidor (\textit{server-side state}). Ao carregar informações de identidade, \textit{tenant} e permissões de forma criptográfica no próprio \textit{token}, o sistema garante que cada requisição seja independente e adequada a ambientes conteinerizados e escaláveis, reduzindo a sobrecarga em bancos de dados de sessão.

Complementando a infraestrutura de desenvolvimento, o uso de ferramentas de conteinerização, como o Docker, reflete uma tendência global indispensável de empacotamento, isolamento e padronização de ambientes de software. Segundo a perspectiva de \citeonline{turnbull2014}, a tecnologia de contêineres permite encapsular a aplicação e todas as suas dependências — como bibliotecas, binários e arquivos de configuração — em uma unidade isolada que compartilha o mesmo núcleo do sistema operacional hospedado.

Essa abordagem garante o princípio da reprodutibilidade, assegurando a consistência absoluta entre os ambientes de desenvolvimento, homologação e produção. Na prática, o isolamento proporcionado pelo Docker mitiga de forma drástica o clássico problema de incompatibilidade operacional conhecido informalmente como ``funciona na minha máquina'', otimizando o pipeline de CI/CD, acelerando o processo de implantação (\textit{deploy}) e oferecendo suporte robusto para a escalabilidade sustentável da aplicação ao longo do tempo.

\section{Computação em Nuvem e CI/CD}

A evolução da computação em nuvem desempenhou um papel central no surgimento de soluções escaláveis, resilientes e acessíveis para organizações de todos os portes. Conforme apontam \citeonline{armbrust}, a computação em nuvem permite que aplicações sejam executadas de forma elástica, proporcionando a alocação dinâmica e automatizada de recursos computacionais de acordo com a demanda instantânea, o que assegura alta disponibilidade e otimização de custos operacionais. 

Dentro desse ecossistema, o modelo baseado em Infraestrutura como Serviço (IaaS), utilizado neste trabalho por meio do provisionamento de uma Instância de Máquina Virtual (VM), destaca-se por oferecer controle total sobre o ambiente de execução e isolamento de recursos. Essa abordagem permite que a aplicação e seus respectivos serviços rodem de maneira dedicada, utilizando a capacidade computacional (CPU, memória e armazenamento) alocada na instância para gerenciar e executar as imagens conteinerizadas do sistema de forma contínua.

Paralelamente, o pipeline de CI/CD, implementado neste projeto por meio do \textit{GitHub Actions}, segue práticas metodológicas amplamente consolidadas pela literatura de \textit{DevOps}. Segundo as diretrizes fundamentais de \citeonline{humble2014}, a automação sistemática dos estágios de compilação (\textit{build}), testes automatizados e implantação (\textit{deploy}) contribui diretamente para entregas de software mais rápidas, confiáveis e frequentes. 

Esse fluxo contínuo reduz drasticamente o risco de falhas humanas em ambiente de produção e viabiliza a evolução incremental do sistema. No cenário desenvolvido, a integração nativa do pipeline com o repositório de artefatos e o processo de entrega automatizado diretamente na instância de VM exemplifica como as arquiteturas modernas combinam de forma sinérgica o versionamento de código, a conteinerização e o \textit{deploy} automatizado, resultando em ciclos de lançamento consistentes, auditáveis e de alta performance.

\subsection{Automação do Pipeline de Implantação e Atualização Contínua}

A arquitetura de um pipeline de implantação moderno exige que a transição entre o código armazenado no repositório e o software em execução na infraestrutura de nuvem ocorra sem intervenção manual. Conforme descrevem \citeonline{kim2016devops}, o objetivo principal de um fluxo de entrega contínua é criar um mecanismo previsível e rotineiro, permitindo que alterações de código sejam integradas e disponibilizadas em ambiente de produção de forma rápida e segura. No modelo estruturado para este projeto, o \textit{GitHub Actions} atua como o motor de orquestração desse fluxo.

O processo inicia-se de forma automatizada a cada modificação na ramificação principal (\textit{main branch}) do repositório da API. O \textit{workflow} gera e publica uma nova imagem \textit{Docker} no GHCR (GitHub Container Registry). Em seguida, o pipeline utiliza SSH e credenciais armazenadas nos \textit{Repository Secrets} do GitHub para acessar a máquina virtual na Google Cloud (Compute Engine).

Uma vez estabelecida a conexão segura, o pipeline realiza o \textit{pull} da imagem publicada no GHCR, atualiza os contêineres via \textit{Docker Compose} e remove imagens antigas. Essa estratégia garante que o ambiente de produção seja atualizado de maneira consistente em poucos minutos. Como destaca \citeonline{humble2014}, a automação desse estágio final elimina os erros comuns de configuração manual em servidores operacionais, consolidando a confiabilidade do software e garantindo que o ERP esteja sempre alinhado com a última versão estável desenvolvida. Destaca-se que o GHCR armazena as imagens de contêiner da API; as mídias da aplicação (arquivos enviados pelos usuários) residem no Google Cloud Storage, conforme detalhado no Capítulo~\ref{cap3MatMet}.

